\documentclass{article}
\usepackage{mymacros,graphicx,setspace} %% changed "doublespace" to "setspace"
\pagestyle{empty}
\setstretch{1.0}

\setlength{\topmargin}{0.0in}
\setlength{\headheight}{0.0in}
\setlength{\headsep}{0.0in}

\setlength{\textheight}{8.2in}
\setlength{\footskip}{0.0in}

\setlength{\textwidth}{6.0in}
\setlength{\oddsidemargin}{0.25in}
\setlength{\evensidemargin}{0.0in}

\begin{document}
\vspace{-3.0in}
\begin{figure}
\centering
\includegraphics[width=2.5in]{nwu}
\end{figure}

%\centerline{\huge {\bf NORTHWESTERN}}
%\vspace{0.1in}
%\centerline{\Large {\bf UNIVERSITY}}
%\vspace{0.3in}
%
\centerline{\Large Department of Electrical Engineering and Computer Science}
\vspace{0.2in}

\centerline{\Large { Technical Report}}
\vspace{0.05in}
\centerline{\Large { NU-EECS-12-04}}
\vspace{0.05in}
\centerline{\Large { October 18, 2012}}
\vspace{0.3in}

% Parser Generator Generator (PGG): A Tool for Rapid Parser Construction via EBNF Grammars
\centerline{\Large {\bf Parser Generator Generator (PGG): A Tool for}}
\centerline{\Large {\bf Rapid Parser Construction via EBNF Grammars}}
\vspace{0.3in}
\centerline{\large {\bf Jedidiah McClurg}}
%  \hspace {0.2in} Goce Trajcevski
\vspace{0.8in}

\centerline{\large {\bf Abstract}}
\vspace{0.15in}
Designing a Domain Specific Language (DSL) can be quite tedious,
with large amounts of time and effort spent building successive
iterations of the lexer and parser as the language evolves
throughout the development process.
Each time a syntactic change is made, both the lexer/parser and
the Abstract Syntax Tree (AST) data structures and manipulation
functions may need to be adjusted.  This process can become very
time-consuming, and can distract from development of the
language's functionality.  Thus, we have built the Parser
Generator Generator (PGG) system to shift some unnecessary
attention away from the frontend implementation details.  PGG accepts a high-level
syntax specification in an expressive variant of Extended Backus Naur Form (EBNF),
and produces corresponding OCaml lexer/parser scripts, as
well as AST data structures and manipulation functions.  The
tool can also emit skeleton code which allows
PGG to be integrated into the project's build sequence,  
enabling projects where the entire lexing/parsing/AST
frontend is specified in a single readable source file,
and allowing DSL developers to quickly get a prototype up and running.

% Designing a Domain Specific Language (DSL) can be quite tedious, with large amounts of time and effort spent building successive iterations of the lexer and parser as the language evolves throughout the development process. Each time a syntactic change is made, both the lexer/parser and the Abstract Syntax Tree (AST) data structures and manipulation functions may need to be adjusted. This process can become very time-consuming, and can distract from development of the language's functionality. Thus, we have built the Parser Generator Generator (PGG) system to shift some unnecessary attention away from the frontend implementation details. PGG accepts a high-level syntax specification in an expressive variant of Extended Backus Naur Form (EBNF), and produces corresponding OCaml lexer/parser scripts, as well as AST data structures and manipulation functions. The tool can also emit skeleton code which allows PGG to be integrated into the project's build sequence, enabling projects where the entire lexing/parsing/AST frontend is specified in a single readable source file, and allowing DSL developers to quickly get a prototype up and running.

\vspace{0.5in}
{\bf Keywords:} Parser, Lexer, AST, EBNF Grammar, OCaml
\end{document}
